\documentclass[12pt]{article}
\usepackage[utf8]{inputenc}
\usepackage[T1]{fontenc}
\usepackage{hyperref}
\title{Yagara Language Resources for K-12 Curriculum Development}
\author{Compiled by Claude Sonnet 5.5 (claude-sonnet-5-5) for Andr\'as Kornai}
\date{}
\begin{document}
\maketitle

\section{Introduction}
Yagara (ISO 639-3: yxg) is a Pama-Nyungan language traditionally spoken in the area that is now Brisbane and Ipswich. There are no fluent native speakers. The Australian English word \emph{yakka} `work' comes from the Yagara verb \emph{yaga} `to work'. A major recent publication supports revitalisation efforts.

\section{Primary Educational Resources}
I found no school primers or textbooks for Yagara. The dictionary below is the most usable base for building classroom materials.

\section{Linguistic and Reference Materials}
\begin{itemize}
\item Karen Sullivan and Glenda Harward-Nalder, \emph{Yagara Dictionary and Salvage Grammar} (ANU Press, 2024). It has over 2,200 dictionary entries with original and standardised spellings, an English word-finder, and all known Yagara texts (sentences, songs and three Bible stories). It is free to download as PDF or EPUB. \url{https://press.anu.edu.au/publications/series/asia-pacific-linguistics/yagara-dictionary-salvage-grammar}
\item The same book in the OAPEN open-access library: \url{https://library.oapen.org/handle/20.500.12657/89785}
\item Language Hat post on the origin of \emph{yakka}: \url{https://languagehat.com/yakka/}
\end{itemize}

\section{Religious and Community Resources}
The book above reproduces three Bible stories in Yagara among the known texts. I found no separate scripture or community organisation in these searches. Co-author Glenda Harward-Nalder is described as a descendant of the Ngugi People of Mulgumpin (Moreton Island), Quandamooka Nation.

\section{Digital and Audio Resources}
I found no audio resources. The book's open-access digital editions are listed above.

\section{Curriculum Development Context}
Yagara is better documented than many extinct languages in this survey. The dictionary, standardised spellings and texts allow word lists, place names and simple phrases to be taught. Use should be coordinated with local Traditional Owners and community language workers.
\end{document}